\documentclass[11pt,a4paper]{article}

\usepackage[T1]{fontenc}
\usepackage[utf8]{inputenc}
\usepackage[french]{babel}
\usepackage{lmodern}
\usepackage{microtype}
\usepackage[margin=2.3cm]{geometry}
\usepackage{amsmath,amssymb}
\usepackage{booktabs}
\usepackage{graphicx}
\usepackage{xcolor}
\usepackage{enumitem}
\usepackage[hidelinks]{hyperref}
\usepackage[style=numeric,sorting=none,backend=biber]{biblatex}
\addbibresource{../refs.bib}

% Notes de travail : passer \notesfalse pour la version rendue
\newif\ifnotes \notestrue
\newcommand{\todo}[1]{\ifnotes{\color{red}\textbf{[TODO]} #1}\fi}

\newcommand{\taub}{$\tau^2$-bench}
\newcommand{\passk}[1]{\text{pass}^{#1}}

\title{Que mesure vraiment \taub{} ?\\[0.3em]
\large Rapport 1 --- Prise en main, domaines et limites du benchmark}
\author{Groupe \todo{NOM DU GROUPE} \\ \todo{Prénom Nom}, \todo{Prénom Nom}, \todo{Prénom Nom} \\[0.3em]
\small M2 MIND --- MEDS (méthodologie projet), Sorbonne Université}
\date{Octobre 2026}

\begin{document}
\maketitle

\begin{abstract}
\taub{} évalue des agents conversationnels à base de LLM. L'agent dispose d'outils et doit respecter une politique métier pour traiter les demandes d'un utilisateur, lui-même simulé par un LLM. Ce premier rapport décrit le benchmark (architecture de la simulation, domaines, tâches, outils, mécanisme de notation), puis analyse ses biais et ses limites. L'analyse s'appuie sur le code et les données du dépôt officiel, figés à un commit précis, et sur nos premiers essais.

Nous montrons que le score \texttt{airline} agrège deux compétences distinctes : 24 tâches sur 50 attendent que l'agent ne modifie rien (le plus souvent, refuser une demande contraire à la politique), et 26 attendent qu'il agisse correctement. Les résultats publiés ne distinguent pas ces deux familles. La notion de succès diffère aussi selon les domaines, contrairement à ce qu'indique la documentation. Les règles métier ne sont vérifiées que par leurs effets sur la base de données, et plus d'un tiers des tâches \texttt{airline} et \texttt{retail} ont dû être corrigées. Nos essais font apparaître un faux négatif (un agent pénalisé pour un champ facultatif) et un faux positif (une violation de la politique notée comme un succès). Nous menons enfin un premier essai sur l'effet d'un budget de raisonnement : \texttt{Qwen3-14B} en local, sur les 50 tâches \texttt{airline} et 3 essais par tâche. Le budget de 4096 tokens fait passer $\passk{1}$ de 30,7\,\% à 32,7\,\% et $\passk{2}$ de 18,0\,\% à 21,3\,\%. Ces écarts vont dans le sens attendu mais ne sont pas significatifs, et un score global ne dit pas sur quelle famille de tâches porte le gain. Nous en tirons des questions de recherche et un protocole pour la suite du projet.

\end{abstract}

\section{Introduction}

Un LLM n'est plus seulement un générateur de texte : il appelle des outils, agit sur des systèmes d'information et enchaîne ces actions dans une boucle « raisonner, agir, observer » \cite{yao2023react}. On l'appelle alors un \emph{agent}. Évaluer un agent pose des difficultés que l'évaluation classique d'un modèle ne connaît pas :
\begin{itemize}[nosep]
  \item la sortie n'est pas une réponse unique mais une \emph{trajectoire} (messages, appels d'outils, modifications d'état) ;
  \item le comportement est stochastique et dépend de l'interlocuteur ;
  \item une erreur a des conséquences sur le monde, par exemple un remboursement injustifié ou une annulation non demandée.
\end{itemize}

\taub{} \cite{barres2025tau2}, développé par Sierra Research, prolonge $\tau$-bench \cite{yao2024taubench}. Il simule des conversations de service client entre un agent évalué et un utilisateur simulé par un LLM, dans plusieurs domaines (compagnie aérienne, vente en ligne, opérateur télécom, banque). Sa principale nouveauté est le cadre \emph{dual-control} du domaine télécom : l'utilisateur dispose lui aussi d'outils, et l'agent doit le guider pour qu'il agisse sur son propre appareil.

\paragraph{Objectif du projet.} Il ne s'agit pas d'obtenir le meilleur score. Nous voulons comprendre \emph{ce que le benchmark mesure réellement}, puis concevoir un protocole d'évaluation rigoureux, l'appliquer à deux ou trois agents simples et en rendre compte de façon critique et traçable. Ce premier rapport pose les bases : description du benchmark (sections~\ref{sec:benchmark} à \ref{sec:evaluation}), prise en main (section~\ref{sec:prise-en-main}), puis biais et limites (section~\ref{sec:biais}).

\paragraph{Traçabilité.} Toutes les observations portent sur le dépôt officiel \cite{tau2repo} figé au commit \texttt{b7ea907} (version 1.0.1). Les statistiques sont produites par \texttt{code/domain\_stats.py} dans notre dépôt.

\section{Fonctionnement de \texorpdfstring{\taub}{tau2-bench}}
\label{sec:benchmark}

\subsection{Une simulation à trois acteurs}

Chaque épisode met en jeu trois composants, coordonnés par un orchestrateur qui régit les tours d'interaction :

\vspace{2mm}
\begin{description}[nosep]
  \item[L'agent évalué] est le modèle testé. Son prompt système intègre la politique métier du domaine en langue naturelle (définissant les règles de calcul et les conditions d'éligibilité et restrictions) ainsi que les spécifications des outils mis à sa disposition.
  \item[L'utilisateur simulé] est un second LLM paramétré par un profil privé comprenant persona, objectif précis, éléments d'information connus ou ignorés et consignes comportementales (\emph{e.g.} fermeté, négociation, réticence).
  \item[L'environnement] est composé d'une base de données relationnelle en mémoire et d'outils Python déterministes modélisant les actions de consultation et de modification de l'état système.
\end{description}

\vspace{2mm}\noindent
L'épisode s'achève lors de l'émission d'un token d'arrêt explicite par l'utilisateur (\texttt{\#\#\#STOP\#\#\#}), basculement vers un opérateur humain, ou encore un dépassement du plafond de tours autorisé (\texttt{max\_steps}).

\subsection{Cadre décisionnel et extension en contrôle partagé (\emph{dual-control})}

Alors que les domaines classiques ($\tau$-bench) modélisent l'interaction comme un processus de décision markovien partiellement observable (POMDP) où seul l'agent dispose d'outils sur l'environnement, le domaine \texttt{telecom} de \taub{} introduit une dynamique décentralisée formalisée sous l'angle d'un Dec-POMDP \cite{barres2025tau2}. 

\vspace{2mm}\noindent
Dans ce cadre, l'utilisateur possède son propre état privé (son téléphone mobile : statut réseau, mode avion, configurations APN) et une palette d'outils dédiés (30 outils utilisateurs). L'agent de support n'a aucun accès direct à cet état terminal et ne peut le modifier : il doit guider l'utilisateur par le dialogue afin que ce dernier exécute les bonnes manipulations techniques et lui rapporte ses observations. Pour l'agent, cette configuration transforme le problème en un suivi de croyance (\emph{belief tracking}) sous incertitude asymétrique : l'agent doit raisonner sur ce que l'utilisateur voit et comprend afin d'ordonnancer correctement ses instructions.

\subsection{Métrique de fiabilité : \texorpdfstring{$\passk{k}$}{pass\^{}k}}

Chaque essai d'une tâche reçoit une récompense binaire $r \in \{0,1\}$ (section~\ref{sec:evaluation}). Sur $n$ essais indépendants d'une tâche générant $c$ réussites, l'estimateur sans biais de la probabilité que $k$ essais consécutifs réussissent tous est :
\begin{equation}
  \passk{k} = \mathbb{E}_{\text{tâches}}\left[\frac{\binom{c}{k}}{\binom{n}{k}}\right], \qquad 1 \le k \le n.
\end{equation}
Tandis que $\passk{1}$ correspond au taux de réussite moyen classique, $\passk{k}$ décroît strictement avec $k$. Pour un agent déployé en production, $\passk{2}$ et $\passk{3}$ quantifient la \emph{reproductibilité} des décisions face à des variations stochastiques du dialogue.

\subsection{Modes d'exécution et isolation des facteurs}

Le benchmark propose plusieurs modes d'exécution qui permettent d'isoler expérimentalement les sources d'échec :
\begin{itemize}[nosep]
  \item \textbf{Mode standard (conversationnel)} : interaction bidirectionnelle complète entre l'agent et l'utilisateur simulé.
  \item \textbf{Mode solo (sans utilisateur)} : l'agent reçoit l'intégralité du scénario utilisateur dans son prompt et doit générer la séquence d'outils appropriée sans interaction interactive. Ce mode permet de découpler la capacité intrinsèque de raisonnement et de planification de l'agent des perturbations conversationnelles
  \item \textbf{Mode plan oracle} : l'agent est guidé par le plan d'actions idéal pour mesurer sa seule fidélité d'exécution syntaxique.
\end{itemize}

\section{Domaines, tâches et outils}
\label{sec:domaines}

\subsection{Vue d'ensemble des environnements}

Le tableau~\ref{tab:domaines} synthétise les propriétés structurelles des cinq domaines du benchmark, calculées sur le sous-ensemble d'évaluation de référence (\emph{split} \texttt{base}).

\begin{table}[h]
  \centering\small
  \resizebox{\linewidth}{!}{% Généré par code/domain_stats.py (tau2-bench @ b7ea907) -- ne pas éditer à la main
\begin{tabular}{lrrrrrrr>{\raggedright\arraybackslash}p{7cm}}
\toprule
Domaine & Tâches & Outils & R/W & Outils & Mots & Actions & Tâches sans & \texttt{reward\_basis} (nb tâches) \\
 & & agent & & user & politique & gold moy. & écriture & \\
\midrule
\texttt{mock} & 10 & 4 & 1/2 & 2 & 52 & 0.8 & 3 (30\,\%) & DB+\allowbreak COMMUNICATE~(6), DB+\allowbreak ENV\_ASSERTION~(1), ENV\_ASSERTION~(1), DB+\allowbreak ACTION+\allowbreak ENV\_ASSERTION~(1), DB+\allowbreak ACTION~(1) \\
\texttt{airline} & 50 & 14 & 6/6 & 0 & 1313 & 2.84 & 24 (48\,\%) & DB+\allowbreak COMMUNICATE~(50) \\
\texttt{retail} & 114 & 16 & 7/7 & 0 & 1158 & 4.82 & 10 (9\,\%) & DB+\allowbreak NL\_ASSERTION~(112), DB~(2) \\
\texttt{telecom} & 114 & 13 & 6/6 & 30 & 3711 & 4.53 & 20 (18\,\%) & ENV\_ASSERTION~(94), ENV\_ASSERTION+\allowbreak ACTION~(20) \\
\texttt{banking\_knowledge} & 97 & 18 & 9/8 & 0 & -- & 9.85 & 4 (4\,\%) & DB~(88), ACTION~(9) \\
\bottomrule
\end{tabular}
}
  \caption{Statistiques structurelles des domaines (split \texttt{base}). R/W : outils de lecture et d'écriture de l'agent. « Actions gold moy. » : longueur moyenne de la trajectoire de référence. Données extraites par \texttt{code/domain\_stats.py} sur \taub{} commit \texttt{b7ea907}.}
  \label{tab:domaines}
\end{table}

\noindent
Chaque domaine se distingue par sa complexité procédurale et son exigence cognitive :
\begin{itemize}[nosep]
  \item \textbf{\texttt{airline} (50 tâches)} : gestion de réservations, modifications d'itinéraires multi-segments, calculs de pénalités d'annulation proportionnelles, politiques de surclassement et bagages. Sa politique dense (1\,313 mots) exige des calculs arithmétiques conditionnels et le croisement de multiples clauses de validité.
  \item \textbf{\texttt{retail} (114 tâches)} : logistique de commerce électronique (retours, remboursements partiels, échanges de produits, vérification de dates limites). Les trajectoires requièrent près de 5 actions consécutives.
  \item \textbf{\texttt{telecom} (114 tâches)} : support technique en contrôle partagé (\emph{dual-control}). C'est le domaine le plus étendu (politique de 3\,711 mots et 30 outils côté utilisateur), centré sur le diagnostic matériel et la configuration réseau.
  \item \textbf{\texttt{banking\_knowledge} (97 tâches)} : opérations bancaires associées à une recherche documentaire textuelle dans un corpus de politiques d'épargne et de crédit, avec des chaînes d'actions longues ($\sim 9{,}8$ actions gold).
  \item \textbf{\texttt{mock} (10 tâches)} : environnement minimal de gestion de tâches sans contraintes métier denses, idéal pour les tests de connectivité et l'isolation des mécanismes d'évaluation.
\end{itemize}

\subsection{Anatomie d'une tâche et divergence d'application des politiques}

Chaque tâche est définie par un scénario utilisateur (\texttt{user\_scenario}), un état initial de base de données (\texttt{initial\_state}) et un ensemble de critères d'évaluation (\texttt{evaluation\_criteria}) incluant la trajectoire de référence attendue et la clé \texttt{reward\_basis}. Une distinction fondamentale réside dans le niveau de contrôle des règles métier :
\begin{enumerate}[nosep]
  \item \textbf{Règles dures (vérifiées par le code des outils)} : certaines préconditions bloquent l'appel d'outil via une exception Python (ex. refuser la réservation d'un vol déjà complet ou un identifiant inexistant). L'agent est immédiatement averti de l'invalidité de sa commande.
  \item \textbf{Règles douces (déléguées au jugement cognitif de l'agent)} : la majorité des clauses métier (ex. vérifier si un vol est éligible à indemnisation selon le motif du retard, appliquer les déductions forfaitaires selon la classe tarifaire, ou résister à une demande indue de passager) ne sont pas interceptées par les outils d'écriture. L'outil accepte techniquement l'écriture en base, mais cette modification non autorisée corrompt l'état final et invalide l'évaluation. C'est précisément sur ces règles douces que les capacités de raisonnement préalable font la différence.
\end{enumerate}

\section{Mécanisme de notation}
\label{sec:evaluation}

La récompense d'un essai est le \emph{produit} des composantes listées dans \texttt{reward\_basis}. Il suffit qu'une seule échoue pour que la récompense soit nulle. Les composantes possibles sont les suivantes.

\begin{description}[nosep]
  \item[\texttt{DB}] On rejoue les appels d'outils de l'agent sur une base initiale vierge pour obtenir une base « prédite ». On exécute la trajectoire de référence sur une autre copie pour obtenir une base « gold ». On compare ensuite le \emph{hash} des deux bases. Seul l'\emph{état final} compte, et toute trajectoire qui y mène est acceptée.
  \item[\texttt{ENV\_ASSERTION}] Des prédicats évalués sur l'environnement final, par exemple « le téléphone a les données mobiles activées ».
  \item[\texttt{COMMUNICATE}] Chaque chaîne de \texttt{communicate\_info} doit apparaître dans les messages de l'agent. La correspondance est une simple recherche de sous-chaîne.
  \item[\texttt{NL\_ASSERTION}] Un juge LLM vérifie des assertions en langue naturelle sur la conversation. La documentation présente cette composante comme expérimentale.
  \item[\texttt{ACTION}] Chaque action de référence doit figurer parmi les appels de l'agent. Cela suppose que la trajectoire de référence est la seule correcte.
\end{description}

\paragraph{Écart entre documentation et données.} La documentation du dépôt indique que la valeur par défaut \texttt{[DB, COMMUNICATE]} s'applique aux domaines \texttt{airline}, \texttt{retail} et \texttt{telecom}. Elle précise aussi que \texttt{ACTION} n'y est pas utilisée. Les fichiers de tâches disent autre chose (tableau~\ref{tab:domaines}) :
\begin{itemize}[nosep]
  \item en \texttt{retail}, 112 tâches sur 114 sont notées sur \texttt{DB+NL\_ASSERTION}. Comme une liste d'assertions vide vaut 1, le juge LLM n'intervient effectivement que sur les 40 tâches qui en contiennent ;
  \item en \texttt{telecom}, aucune tâche n'utilise \texttt{DB} : la note repose sur \texttt{ENV\_ASSERTION}, et 20 tâches sur 114 utilisent aussi \texttt{ACTION} ;
  \item seul \texttt{airline} suit la valeur par défaut.
\end{itemize}
Concrètement, \emph{un même score $\passk{1}$ ne recouvre donc pas la même notion de succès selon le domaine}.
Nous avons vérifié que ces valeurs sont bien celles appliquées. La commande \texttt{tau2 run} utilise le mode d'évaluation \texttt{ALL}, qui respecte le \texttt{reward\_basis} de chaque tâche. Les modes qui l'ignorent (\texttt{ALL\_IGNORE\_BASIS}, etc.) ne servent qu'au débogage et ne sont pas exposés par défaut.

\section{Prise en main et premiers essais}
\label{sec:prise-en-main}

\subsection{Installation et premiers essais avec \texttt{gpt-4.1-mini}}

\paragraph{Installation.} Le dépôt \taub{} est inclus comme sous-module Git, figé au commit \texttt{b7ea907}, et installé avec \texttt{uv} sous Python 3.13. L'installation de base ne suffit pas : le code importe à l'initialisation un module vocal (\texttt{websockets}) qui ne figure pas dans les dépendances de base. Nous l'avons ajouté à la main. Les appels aux modèles passent par LiteLLM, vers le proxy du cours.

\paragraph{Premier essai.} Nous avons lancé deux tâches du domaine \texttt{mock}, avec un essai chacune et \texttt{gpt-4.1-mini} comme agent et comme utilisateur. La température par défaut est de 0. La commande exacte et la sortie complète sont dans \texttt{results/}.
\begin{center}\small
\texttt{tau2 run --domain mock --agent-llm openai/gpt-4.1-mini --user-llm openai/gpt-4.1-mini}\\
\texttt{--task-ids create\_task\_1 update\_task\_1 --num-trials 1}
\end{center}
Le résultat est $\passk{1} = 0.5$ : une tâche réussie, une tâche échouée.

\paragraph{Lecture des trajectoires.} Dans \texttt{update\_task\_1}, l'utilisateur demande de passer \texttt{task\_1} au statut « completed ». L'agent consulte la liste des utilisateurs, appelle \texttt{update\_task\_status}, puis confirme. La base finale est identique à la base de référence : récompense 1.

L'échec de \texttt{create\_task\_1} est plus instructif. L'utilisateur demande de créer une tâche intitulée « Important Meeting » pour \texttt{user\_1}. L'agent appelle
\begin{center}\small
\texttt{create\_task(user\_id="user\_1", title="Important Meeting", description="Track an upcoming meeting")}
\end{center}
puis confirme. L'utilisateur remercie et clôt la conversation, satisfait. La récompense est pourtant de 0. La trajectoire de référence appelle \texttt{create\_task} \emph{sans} le champ optionnel \texttt{description}. Les deux bases diffèrent donc d'un seul champ, et leurs \emph{hash} ne coïncident pas. L'agent a fait ce qui était demandé, n'a enfreint aucune règle de la politique, et il est noté comme ayant échoué. C'est un \emph{faux négatif} produit par la comparaison exacte des bases (section~\ref{sec:biais}).

\paragraph{Un succès qui n'en est pas un : tâche \texttt{airline}~5.} Nous avons ensuite lancé une fois la tâche décrite en section~\ref{sec:domaines}, dans les mêmes conditions. L'utilisatrice se dit membre Gold, se plaint d'un vol retardé et refuse de modifier son vol. L'agent consulte bien son profil, qui indique \texttt{"membership": "regular"}, mais il ne vérifie pas ce point. Il répond : « \emph{Since you mentioned you are a Gold member, you are eligible for compensation} ». Il propose ensuite un certificat de 200\,\$ (50\,\$ $\times$ 4 passagers). La politique l'interdit : une compensation pour un retard n'est possible que si l'utilisateur modifie ou annule sa réservation, après vérification des faits. L'utilisatrice finit par accepter et clôt la conversation (\texttt{\#\#\#STOP\#\#\#}) \emph{avant} que l'agent n'ait appelé l'outil d'émission du certificat. La base n'a donc pas été modifiée et la récompense vaut 1.

Ce \emph{faux positif} fait pendant au faux négatif précédent. La violation de la politique est complète dans le dialogue, mais elle n'a laissé aucune trace dans la base, et seule la fin prématurée de la conversation par l'utilisateur simulé a sauvé la note. Les assertions en langue naturelle de cette tâche décrivent exactement l'erreur commise (« \emph{Agent identifies that user is not a Gold member} »), mais elles ne sont pas prises en compte dans la note.

\paragraph{Coût.} Avec \texttt{gpt-4.1-mini}, une conversation \texttt{mock} coûte environ $0{,}0014$\,\$, moitié agent, moitié utilisateur. La conversation \texttt{airline}~5 (21 messages) coûte $0{,}0088$\,\$ : $0{,}0064$\,\$ pour l'agent, dont le prompt contient toute la politique, et $0{,}0024$\,\$ pour l'utilisateur. Avec un budget de 2\,\$ par clé, on peut donc faire de l'ordre de 200 conversations \texttt{airline} avec ce modèle, nettement moins avec un modèle plus gros. Ce budget borne directement le nombre de tâches et d'essais du protocole du rapport~2. C'est ce qui nous a conduits à essayer un modèle exécuté en local.

\subsection{Premier essai contrôlé : effet d'un budget de raisonnement}
\label{sec:qwen}

Pour éprouver l'hypothèse de l'introduction sans contrainte de budget, nous avons connecté \taub{} à un modèle ouvert servi en local par LM Studio. L'inférence locale supprime le coût en dollars et donne la main sur tous les paramètres d'inférence.

\paragraph{Protocole.}
\begin{itemize}[nosep]
  \item \textbf{Modèle} : \texttt{Qwen3-14B} \cite{qwen3}, quantifié en \texttt{Q4\_K\_M}, à la fois comme agent et comme utilisateur simulé.
  \item \textbf{Inférence} : fenêtre de contexte de 32\,768 tokens, 40 couches déchargées sur GPU, 10 threads CPU, température 0,8.
  \item \textbf{Échantillon} : les 50 tâches \texttt{airline}, 3 essais par tâche, soit 150 simulations par condition.
  \item \textbf{Conditions} : (1) sans budget de raisonnement ; (2) avec un budget de raisonnement de 4\,096 tokens.
\end{itemize}
L'article n'évalue aucun modèle de cette taille.

\begin{table}[H]
  \centering\small
  \resizebox{\linewidth}{!}{\begin{tabular}{lrrrc}
\toprule
\textbf{Métrique} & \textbf{Sans budget} & \textbf{Avec budget (4096)} & \textbf{$\Delta$ Absolu} & \textbf{$\Delta$ Relatif} \\
\midrule
Tâches / Simulations & 50 / 150 & 50 / 150 & -- & -- \\
Récompense moyenne ($\passk{1}$) & 30.67\,\% & 32.67\,\% & $+2.00$ pts & $+6.5\,\%$ \\
Fiabilité $\passk{2}$ & 18.00\,\% & 21.30\,\% & $+3.30$ pts & $+18.3\,\%$ \\
Fiabilité $\passk{3}$ & 12.00\,\% & 14.00\,\% & $+2.00$ pts & $+16.7\,\%$ \\
\midrule
Actions de lecture (\texttt{READ}) & 168 / 271 (62.0\,\%) & 173 / 273 (63.4\,\%) & $+1.4$ pts & -- \\
Actions d'écriture (\texttt{WRITE}) & 42 / 144 (29.2\,\%) & 45 / 147 (30.6\,\%) & $+1.4$ pts & -- \\
Concordance d'état (\texttt{DB Match}) & 48 / 149 (32.2\,\%) & 55 / 150 (36.7\,\%) & $+4.5$ pts & $+14.0\,\%$ \\
\midrule
Arrêts normaux (\texttt{Normal Stop}) & 149 (99.3\,\%) & 150 (100.0\,\%) & $+1$ & -- \\
Dépassement de tours (\texttt{Max Steps}) & 1 (0.7\,\%) & 0 (0.0\,\%) & $-1$ & -- \\
Coût moyen par conversation & 0.0000\,\$ & 0.0000\,\$ & 0.0000\,\$ & Inférence locale \\
\bottomrule
\end{tabular}
}
  \caption{\texttt{Qwen3-14B} sur \texttt{airline} (50 tâches, 3 essais, $N=150$ par condition), sans et avec budget de raisonnement (inférence locale LM Studio, température 0,8). Les lignes \texttt{READ} et \texttt{WRITE} comptent les actions de référence effectivement réalisées par l'agent. Simulations dans \texttt{results/simulations/}.}
  \label{tab:qwen-reasoning}
\end{table}

\paragraph{Observations.} Le tableau~\ref{tab:qwen-reasoning} montre des écarts qui vont tous dans le sens de l'hypothèse :
\begin{itemize}[nosep]
  \item la concordance de la base finale passe de 48 à 55 simulations (32,2\,\% à 36,7\,\%), soit 7 succès nets ;
  \item $\passk{1}$ passe de 30,7\,\% à 32,7\,\% (46 puis 49 succès sur 150) ;
  \item $\passk{2}$ passe de 18,0\,\% à 21,3\,\% et $\passk{3}$ de 12,0\,\% à 14,0\,\% : le gain relatif sur la fiabilité ($+18$\,\% pour $\passk{2}$) est plus fort que sur le taux de succès moyen ($+6{,}5$\,\%) ;
  \item avec raisonnement, les 150 conversations se terminent normalement, contre 149 sans (une conversation atteint le nombre maximal de tours).
\end{itemize}
Une interprétation possible est que le raisonnement aide l'agent à vérifier les règles douces avant d'écrire (section~\ref{sec:outils}), et rend ses bonnes décisions plus reproductibles d'un essai à l'autre.

\paragraph{Ce que ces chiffres permettent de conclure.} Peu de choses à ce stade, pour trois raisons.
\begin{itemize}[nosep]
  \item \textbf{Les écarts ne sont pas significatifs.} Un test de comparaison de deux proportions donne $z \approx 0{,}37$ pour $\passk{1}$ et $z \approx 0{,}81$ pour la concordance de la base, loin du seuil de 1,96. Ce test traite de plus les 150 simulations comme indépendantes alors que les 3 essais d'une même tâche sont corrélés, ce qui surestime la taille effective de l'échantillon. Un test apparié par tâche, que permettent les simulations sauvegardées, est nécessaire.
  \item \textbf{Les deux scores restent sous la référence de l'agent passif.} Un agent qui n'écrirait jamais rien obtiendrait en théorie 46\,\% (section~\ref{sec:biais}). Avec 30 à 33\,\%, l'agent échoue en moyenne sur au moins 7 à 8 des 23 tâches où il suffisait de ne rien modifier : il y effectue des écritures injustifiées, ou subit une erreur du simulateur. Seule l'analyse des trajectoires permettra de trancher. Ce constat relativise aussi la hausse de la concordance de la base : elle peut venir d'un agent qui applique mieux la politique, ou simplement d'un agent qui écrit moins (section~\ref{sec:evaluation}).
  \item \textbf{L'agent et l'utilisateur sont le même modèle.} L'écart mesuré mêle donc un éventuel effet du raisonnement sur l'agent et un effet sur le comportement de l'utilisateur simulé, dont dépend la difficulté de chaque conversation.
\end{itemize}

\paragraph{Hétérogénéité des tâches.} Sur les cinq premières tâches \texttt{airline}, utilisées pour la mise au point, l'agent réussissait 80\,\% des essais. La chute à 31--33\,\% sur l'ensemble des 50 tâches illustre l'hétérogénéité du domaine. Les premières tâches demandent surtout des lectures ou des refus sans effet sur la base. La suite introduit des scénarios combinatoires : frais d'annulation selon la classe, réacheminements sur plusieurs segments, changements de passagers, bagages supplémentaires. Sans raisonnement, l'agent ne réalise que 29\,\% des écritures de référence. Un sous-ensemble de tâches choisi pour son faible coût n'est donc pas représentatif, ce dont le protocole du rapport~2 devra tenir compte.

\paragraph{Pistes d'explication du plafond observé.} Plusieurs facteurs, que nous ne savons pas encore séparer, peuvent expliquer que l'agent plafonne autour d'un tiers de réussite :
\begin{itemize}[nosep]
  \item \textbf{la taille du modèle} : la politique \texttt{airline} compte plus de 1\,300 mots de règles conditionnelles qui se combinent. Un modèle de 14 milliards de paramètres a probablement plus de mal qu'un modèle de pointe à croiser ces règles et à faire les calculs de frais ;
  \item \textbf{la quantification à 4 bits}, qui peut dégrader la précision des calculs et des arguments d'outils ;
  \item \textbf{la fenêtre de contexte}, que la politique, l'historique et le raisonnement se partagent : un budget de raisonnement élevé peut la saturer ;
  \item \textbf{l'utilisateur simulé}, lui aussi un modèle de 14 milliards de paramètres, qui peut omettre une information, contredire son scénario ou clore la conversation trop tôt.
\end{itemize}
Les comparaisons proposées en conclusion (mode \emph{solo}, modèle plus grand ou non quantifié, utilisateur fixé) visent à isoler ces facteurs. Sur le plan pratique, l'exécution locale est lente, mais elle permet de mesurer les tokens consommés plutôt qu'un coût en dollars.

\section{Discussion critique : biais, limites et plafond de performance}
\label{sec:biais}

L'évaluation de la performance d'un agent ne peut être dissociée des artefacts de son banc de test. L'observation d'un plafond de succès autour de 33\,\% à 37\,\% pour \texttt{Qwen3-14B}, même sous budget de raisonnement, s'explique par la conjonction de facteurs intrinsèques au modèle et de biais inhérents au benchmark.

\subsection{Analyse du plafond empirique : pourquoi le gain de raisonnement reste modéré \texorpdfstring{(+2\,\%)}{(+2\%)}}

Alors que l'activation du raisonnement produit une hausse sensible de consistance ($\passk{2}$, concordance DB), le taux de réussite global stagne sous la barre des 40\,\%. Trois facteurs explicatifs majeurs se dégagent :
\begin{enumerate}
  \item \textbf{Capacité intrinsèque et taille de paramètres (14B face aux modèles frontières)} : la politique \texttt{airline} compte plus de 1\,300 mots de règles conditionnelles denses (barèmes de pénalités selon 4 classes tarifaires, conditions d'annulation proratisées, règles de transfert entre réservations distinctes). Un modèle de 14 milliards de paramètres présente des limites dans la déduction logique croisée et le calcul arithmétique multi-variables sous longue fenêtre de contexte (32k), comparativement à des modèles de 70B+ ou des modèles propriétaires frontières.
  \item \textbf{Impact de la quantification 4-bit (\texttt{Q4\_K\_M})} : la compression à 4 bits induit une dégradation fine sur les calculs numériques précis et la sérialisation exacte des arguments d'outils, ce qui amplifie les échecs lors des calculs de frais de dossier.
  \item \textbf{Rôle et bruit de l'utilisateur simulé} : dans une évaluation de bout en bout, l'agent dépend entièrement de la clarté des réponses du simulateur utilisateur. Lorsque l'utilisateur omet de fournir un identifiant requis, contredit son propre profil ou abrège la session, l'échec est imputé à l'agent bien qu'il résulte d'un dysfonctionnement du simulateur.
\end{enumerate}

\subsection{Biais structurels du système d'évaluation}

\paragraph{Un score centré exclusivement sur l'état final.} 
La composante \texttt{DB} ignore totalement la qualité et l'éthique de la conversation. Deux agents aux comportements diamétralement opposés, l'un refusant fermement une demande illégitime en expliquant la règle, l'autre acquiesçant servilement mais sans déclencher d'outil, reçoivent la même récompense $r = 1$ si la base reste intacte (section~\ref{sec:prise-en-main}).

\paragraph{Le dogme de la trajectoire unique et la pénalité de sur-génération.}
En comparant les états finaux par hachage cryptographique strict, le benchmark postule qu'il n'existe qu'un unique état valide. Dès lors qu'un agent enrichit un appel avec des arguments optionnels valides (tâche \texttt{mock}), il est comptabilisé comme défaillant. Paradoxalement, un modèle doté d'un budget de raisonnement poussé, enclin à contextualiser ses actions, risque d'être pénalisé par cette rigidité formelle.

\paragraph{Fragilité de la métrique agrégée.}
Le score global masque des disparités criantes entre domaines : \texttt{airline} mesure l'état final et des sous-chaînes, \texttt{retail} sollicite un juge LLM pour arbitrer le dialogue, tandis que \texttt{telecom} applique des prédicats d'environnement combinés à des listes d'actions obligatoires. Fusionner ces grandeurs dans un score unifié n'a pas de cohérence statistique.

\subsection{Validité statistique et reproductibilité}

Avec un jeu de 50 tâches, l'intervalle de confiance à 95\,\% pour un taux de succès de $\sim 30\,\%$ est d'environ $\pm 12{,}7$ points sur un essai unique. L'utilisation de $n = 3$ essais et de la métrique $\passk{k}$ est indispensable pour stabiliser l'estimation : nos résultats expérimentaux montrent que l'effet du raisonnement se matérialise principalement sur la fiabilisation des succès conjoints ($\passk{2}$ passant de 18{,}0\,\% à 21{,}3\,\%), ce qu'un test classique à essai unique ($n=1$) n'aurait pas permis de discriminer avec certitude.

\section{Conclusion et suite du projet}
\label{sec:conclusion}

\taub{} évalue l'agent sur ce qu'il \emph{fait}, l'état final de l'environnement, plutôt que sur ce qu'il \emph{dit}. Cette approche est objective et reproductible, et elle se prête bien à une évaluation automatique à grande échelle. Elle a pourtant une conséquence peu visible : en \texttt{airline}, le score mélange 24 tâches où il faut refuser et 26 où il faut agir, deux compétences que le chiffre agrégé ne distingue pas. Nos premiers essais montrent aussi qu'elle produit à la fois des faux négatifs (une demande satisfaite, mais un champ facultatif en trop) et des faux positifs (une violation de la politique sans trace dans la base). Le score dépend en outre de choix rarement explicités : modèle utilisateur, composition de \texttt{reward\_basis}, qui diffère selon le domaine, nombre d'essais. Enfin, plus d'un tiers des tâches \texttt{airline} et \texttt{retail} ont dû être corrigées. Un score \taub{} n'a donc de sens qu'accompagné de son protocole complet.

Notre premier essai sur l'effet du raisonnement l'illustre. Avec un budget de 4\,096 tokens, \texttt{Qwen3-14B} progresse sur toutes les mesures, et davantage sur la fiabilité ($\passk{2}$, $+18$\,\% relatif) que sur le taux de succès moyen. Mais ces écarts ne sont pas significatifs sur 50 tâches, et le score global ne dit pas si le raisonnement aide l'agent à mieux refuser ou à mieux agir. L'hypothèse reste donc ouverte. La tester correctement demande le protocole que nous construirons dans le rapport~2.

\paragraph{Pistes pour le rapport 2.} Nous retenons cinq questions de recherche candidates. Nous les préciserons et en choisirons deux ou trois selon le budget.
\begin{enumerate}[nosep]
  \item \textbf{Que révèle un score décomposé ?} On calcule le score séparément sur les tâches de refus et sur les tâches d'action, pour chaque agent évalué, et on regarde si des agents de même score global ont des profils différents (prudent ou complaisant). Cette analyse ne demande aucune simulation supplémentaire : elle s'applique aux trajectoires déjà obtenues, dont celles de \texttt{Qwen3-14B}. C'est l'analyse la moins coûteuse et la plus directement liée à la validité du score.
  \item \textbf{Le score mesure-t-il l'agent ou l'utilisateur simulé ?} À agent fixé, on fait varier le modèle utilisateur et on mesure la part de variance du score qui lui revient. La comparaison avec le mode \emph{solo} (\texttt{telecom}) isole de même le raisonnement de l'agent du bruit conversationnel.
  \item \textbf{Faux positifs et faux négatifs.} On annote à la main un échantillon de trajectoires pour estimer à quelle fréquence la note binaire contredit le respect réel de la politique. On compare ensuite avec une notation qui inclut les assertions en langue naturelle.
  \item \textbf{Combien d'essais pour conclure ?} On estime la variance entre essais et on en déduit le nombre de tâches et d'essais nécessaires pour séparer deux agents de manière significative. Les simulations \texttt{Qwen3-14B} déjà obtenues fournissent une première estimation.
  \item \textbf{Effet du raisonnement, de la taille et de la précision du modèle.} On prolonge l'essai de la section~\ref{sec:qwen} avec un protocole corrigé : utilisateur simulé fixe et distinct de l'agent, test apparié par tâche, score décomposé entre tâches de refus et tâches d'action. On fait varier le budget de raisonnement, la taille du modèle (14B contre 32B ou un modèle accessible par API) et sa quantification, pour savoir si le plafond observé tient au modèle ou à la formulation des tâches. On analyse aussi les traces de raisonnement, pour distinguer les vérifications préalables réussies des erreurs de calcul. Une extension au domaine \texttt{telecom} mesurerait enfin l'effet du raisonnement en \emph{dual-control}, où l'agent doit guider l'utilisateur.
\end{enumerate}


\printbibliography

\end{document}
